\section*{第 \ref{ch:b2:microbial-metabolism} 章——微生物代谢与生物膜}

\begin{solution}{exo:b2:microbial-metabolism:1}
蓝细菌：光能无机自养型（光、水、$\mathrm{CO_2}$）。
\emph{Nitrosomonas}：化能无机自养型（铵氧化，
$\mathrm{CO_2}$）。\emph{E.~coli} 以葡萄糖为底物：化能有机异养型。
紫色非硫细菌在光照下以琥珀酸盐为底物：
光能有机异养型。\emph{Thiobacillus} 在黑暗中以硫化物为底物：
化能无机自养型。
\end{solution}

\begin{solution}{exo:b2:microbial-metabolism:2}
氧气（表层数毫米），硝酸盐（稍下方，氧气
已耗尽之处），锰和铁的氧化物（棕色向灰色的转变处），
硫酸盐（黑色硫化物层，海洋沉积物中较厚），
产甲烷所用的 $\mathrm{CO_2}$（最深处，或在硫酸盐贫乏的淡水中紧接
表层之下）。
\end{solution}

\begin{solution}{exo:b2:microbial-metabolism:3}
$5\,\mathrm{C_6H_{12}O_6} + 24\,\mathrm{NO_3^-} + 24\,\mathrm{H^+}
\to 30\,\mathrm{CO_2} + 12\,\mathrm{N_2} + 42\,\mathrm{H_2O}$：4.8
个硝酸根离子，即每个葡萄糖分子（每个硝酸根接受 5 个电子，
每个葡萄糖释放 24 个电子）。
\end{solution}

\begin{solution}{exo:b2:microbial-metabolism:4}
一种附着于表面的微生物群落，包埋于自身分泌的
由多糖、蛋白质和 DNA 构成的基质中。阶段：附着、
微菌落形成与基质分泌、带菌塔与
通道的成熟、扩散。牙菌斑与龋齿；导管和植入物的感染
（包括囊性纤维化患者的肺部）；管道与船体的污损
和腐蚀——或者，有益地，
用作水处理的滴滤池。
\end{solution}

\begin{solution}{exo:b2:microbial-metabolism:5}
八个电子由 \qty{-0.42}{V} 降至 \qty{-0.24}{V}：
$\Delta G^{\circ\prime} = -8\times 96.5\times 0.18 =
\qty{-139}{kJ/mol}$（每摩尔甲烷；查表值为 \qty{-131}{kJ/mol}），即
每个 $\mathrm{H_2}$ 约 \qty{35}{kJ}：每个氢至多 0.7 个 ATP，
实际约为其一半。每个底物分子的能量如此
之少，以致产甲烷菌生长缓慢（\omterm{def:b2:unicellular-diversity:division}{世代时间}为数小时至
数天），几乎将全部底物转化为气体，而非
细胞。
\end{solution}

\begin{solution}{exo:b2:microbial-metabolism:6}
$\mu = 1.2\,S/(5 + S)$：\qty{0.20}{h^{-1}}（$T = \ln 2/\mu =
\qty{3.5}{h}$），\qty{0.60}{h^{-1}}（\qty{1.2}{h}），\qty{1.09}{h^{-1}}
（\qty{0.64}{h}）。达到九成时：$S = 9K_s = \qty{45}{mg/L}$。
\end{solution}

\begin{solution}{exo:b2:microbial-metabolism:7}
$S^* = 0.05\times 0.4/(0.8 - 0.4) = \qty{0.05}{g/L}$; $X^* =
0.4\times(5 - 0.05) = \qty{1.98}{g/L}$；输出 $DX^* = 0.4\times
1.98 = \qty{0.79}{g/L/h}$；洗出 $D_c = 0.8\times 5/5.05 =
\qty{0.79}{h^{-1}}$。
\end{solution}

\begin{solution}{exo:b2:microbial-metabolism:8}
捕获量：$0.25\times 275 = \qty{69}{kJ}$，即每个铵离子
产生 1.4 个 ATP；$12/1.4 = 8.7$，约每个碳需九个铵离子——
还须计入由亚硝酸盐逆向驱动电子以产生 NADH 的代价，则
需数倍之多。异养生物所得的碳已
被还原且有组织，每个葡萄糖可得三十个 ATP；硝化细菌必须
从贫乏的供体中获取能量与还原力，并
由 $\mathrm{CO_2}$ 构建全部碳，因而它只将所处理底物的
极小部分转化为细胞。
\end{solution}

\begin{solution}{exo:b2:microbial-metabolism:9}
$L_p = \sqrt{2\times 1.2\times 10^{-9}\times 0.2/0.5} =
\sqrt{9.6\times 10^{-10}} = \qty{31}{\micro m}$。加倍 $c_0$
使 $L_p$ 乘以 $\sqrt 2$：\qty{44}{\micro m}；将 $q$ 增至四倍
则使其减半：\qty{15}{\micro m}。
\end{solution}

\begin{solution}{exo:b2:microbial-metabolism:10}
自上而下：含藻类与蓝细菌的水体（产氧
光合作用）；铁氧化菌形成的锈色带，此处 $\mathrm{Fe^{2+}}$
与氧相遇；无色硫氧化菌形成的白色薄膜，
（化能\omterm{def:b2:microbial-metabolism:trophic}{无机营养生物}），此处硫化物与氧相遇；紫色带
由紫硫细菌构成（以硫化物为电子供体的\omterm{def:b2:microbial-metabolism:anoxygenic}{不产氧光合作用}）；
绿色带为绿硫细菌（同样过程，能耐受更多
硫化物和更弱的光照）；黑色淤泥，由\omterm{def:b2:microbial-metabolism:fermentation}{发酵}菌与\omterm{def:b2:microbial-metabolism:anaerobic}{硫酸盐还原}菌构成。
(a) 紫硫细菌既需上方的光，又需下方的硫化物，
因此它们位于两种梯度相交的界面处。
(b) \omterm{def:b2:microbial-metabolism:anaerobic}{硫酸盐还原}菌产生 $\mathrm{H_2S}$，使铁以
黑色 FeS 的形式沉淀。(c) 就物质而言封闭：硫、碳与氮在
氧化态与还原态之间循环而不外流；就能量而言开放：
光进入，热散失，若无光，每一循环
均将停止。
\end{solution}

\begin{solution}{exo:b2:microbial-metabolism:11}
产生速率 $pn$，损失 $kA$：$A^* = pn/k$。$A_c = 10^{-8}\times
6.02\times 10^{23} = 6.0\times 10^{15}$ 个分子每升；$n_c =
kA_c/p = 5\times 10^{-4}\times 6.0\times 10^{15}/5 = 6\times 10^{11}$
每升，即每毫升 $6\times 10^{8}$ 个。稠密培养
（$10^{9}$）高于阈值，海水（$10^{5}$）低于阈值达四个
数量级。\omterm{def:b2:microbial-metabolism:biofilm}{生物膜}中细胞数达 $10^{11}$ 个每
毫升基质，而基质减缓了信号的逸出
（实际上降低了 $k$），故仅数千个细胞的微菌落
位于纳升体积内即可达到\omterm{prop:b2:microbial-metabolism:quorum}{群体感应}阈值，而同样数量的细胞分散于一
升之中则永远不会达到。
\end{solution}

\begin{solution}{exo:b2:microbial-metabolism:12}
固氮（只有\omterm{def:b2:unicellular-diversity:prokaryote}{原核生物}具有固氮酶）：停止后，生物圈的
化合态氮将经反硝化而流失，
生产力会在数十年内崩溃。
\omterm{prop:b2:microbial-metabolism:lithotrophy}{硝化作用}与\omterm{def:b2:microbial-metabolism:anaerobic}{反硝化作用}（仅\omterm{def:b2:unicellular-diversity:prokaryote}{原核生物}）：停止后，
铵将积累、硝酸盐消失，氮将
不再回到空气中。\omterm{def:b2:microbial-metabolism:anaerobic}{产甲烷作用}与无氧呼吸
（仅\omterm{def:b2:unicellular-diversity:prokaryote}{原核生物}）：停止后，缺氧沉积物、瘤胃与沼泽中的有机物
将无法矿化而积累，碳将
退出循环。同样，产氧光合作用本是细菌的
发明，无此则无可供呼吸的氧气。真核生物
是\omterm{def:b2:unicellular-diversity:prokaryote}{原核生物}星球上晚近、化学上范围较窄的
补充。
\end{solution}

\begin{solution}{pb:b2:microbial-metabolism:1}
\textbf{1.} $-24\times 96.5\times (0.82 + 0.43) = \qty{-2900}{kJ/mol}$.
\textbf{2.} 硝酸盐：$-24\times 96.5\times 1.17 = \qty{-2700}{kJ/mol}$；
硫酸盐：$-24\times 96.5\times 0.21 = \qty{-490}{kJ/mol}$。
\textbf{3.} 氧的产量更多，因此只要有氧，好氧菌就会在争夺有机物时胜过
反硝化细菌，氧还会抑制
硝酸盐还原酶；\omterm{def:b2:microbial-metabolism:anaerobic}{硫酸盐还原}需要缺氧\emph{且}
无硝酸盐，因此其异味说明池中曝气不足，
硝酸盐已耗尽——而 $\mathrm{H_2S}$ 有毒，并会腐蚀
混凝土。
\textbf{4.} $0.4\times 2900/50 = 23$ ATP；$0.4\times 2700/50 = 22$；
$0.4\times 490/50 = 4$.
\textbf{5.} 约多出六倍（$23/4$）。
\textbf{6.} 每个葡萄糖产生两个 ATP，意味着产量约为 \num{0.1}，因此
\qty{90}{\%}的底物能量仍留在产物中（酸、
醇、$\mathrm{H_2}$），这些产物由产甲烷菌转化为甲烷，
其获益同样很少；能量少，生物量也少，并且
碳以气体形式离开系统。
\textbf{7.} 每天 COD 为 $10^4\times 0.3 = \qty{3000}{kg}$。
\textbf{8.} 生物量对应的 COD 为 $0.4\times 3000 = \qty{1200}{kg}$；
剩余 \qty{1800}{kg} 的 COD 被氧化，需要 \qty{1800}{kg}
的氧气 $\mathrm{O_2}$。
\textbf{9.} 每天 $1800/2 = \qty{900}{kWh}$。
\textbf{10.} 氮为 $10^4\times 0.04 = \qty{400}{kg}$；
需氧气 $4.57\times 400 = \qty{1830}{kg}$（$\mathrm{O_2}$）。
\textbf{11.} 稳态时 $\mu = 1/\theta \le \mu_{\max}$，
洗出发生于 $\theta = 1/\mu_{\max} = \qty{2}{d}$ 以下。
\textbf{12.} $D = \qty{0.1}{d^{-1}}$: $S^* = 1\times 0.1/(0.5 - 0.1)
= \qty{0.25}{g/m^3}$.
\textbf{13.} $S^* = 0.1/(0.25 - 0.1) = \qty{0.67}{g/m^3}$，且
洗出极限升至 \qty{4}{d}；操作者须延长
污泥龄（少排污泥）：在 $\theta = \qty{20}{d}$ 时，$S^* =
0.05/0.2 = \qty{0.25}{g/m^3}$，仍为原值。
\textbf{14.} 每天 COD 为 $2.86\times 400 = \qty{1144}{kg}$。
\textbf{15.} 出水中铵为 $0.25\times 10^4 = \qty{2.5}{kg/d}$，
占负荷的 \qty{0.6}{\%}；其余 \qty{397}{kg} 的氮，
以 $\mathrm{N_2}$ 形式离开，占比 \qty{99}{\%}。
\textbf{16.} 节省 \qty{1144}{kg} 的 $\mathrm{O_2}$。总需求：
每天 $(1800 - 1144) + 1830 = \qty{2490}{kg}$ 的 $\mathrm{O_2}$。
\textbf{17.} $1200\times 0.6\times 0.35 = \qty{252}{m^3}$ 的甲烷，
每天。
\textbf{18.} $252\times 36 = \qty{9070}{MJ} = \qty{2520}{kWh}$；
每天电量 $0.35\times 2520 = \qty{880}{kWh}$。
\textbf{19.} 曝气 $2490/2 = \qty{1245}{kWh}$，每天：甲烷
可满足其 $880/1245 = \qty{71}{\%}$。
\textbf{20.} 由 $D_e\,c'' = q$，得 $c = c_0 + ax + qx^2/2D_e$；代入
$c(L_p) = 0$ 与 $c'(L_p) = 0$，得 $a = -qL_p/D_e$，且 $c =
(q/2D_e)(L_p - x)^2$，并有 $L_p^2 = 2D_ec_0/q$。
\textbf{21.} $L_p = \sqrt{2\times 1.5\times 10^{-9}\times 0.2/0.3} =
\sqrt{2\times 10^{-9}} = \qty{45}{\micro m}$.
\textbf{22.} 好氧部分占 $45/300 = \qty{15}{\%}$。下层的反硝化细菌
将从硝化表层向下扩散的硝酸盐
还原为 $\mathrm{N_2}$，利用向内扩散的有机物。
\textbf{23.} 在同一\omterm{def:b2:microbial-metabolism:biofilm}{生物膜}内，氧梯度形成好氧层
（\omterm{prop:b2:microbial-metabolism:lithotrophy}{硝化作用}）覆于缺氧层（\omterm{def:b2:microbial-metabolism:anaerobic}{反硝化作用}）之上，两层
相距数十微米；曝气池处处有氧，因此
反硝化需要单独的缺氧池。
\textbf{24.} 氯在到达深层细胞之前，已与基质
聚合物反应并被消耗；深层细胞
处于饥饿、生长缓慢或休眠状态，敏感性低得多；
仅表层死亡，\omterm{def:b2:microbial-metabolism:biofilm}{生物膜}从下方重新生长。
\textbf{25.} 每天氧需求 \qty{2490}{kg}；每天甲烷
\qty{252}{m^3}；甲烷供应曝气电量的 \qty{71}{\%}
（\qty{880}{kWh}，占 \qty{1245}{kWh}）。
\end{solution}
